% These summaries are generated from the reviewed 232-row v4 campaign table.
\providecommand{\VFourLocalPZeroMinimum}{$1.7636\times10^{-4}$}
\providecommand{\VFourLocalPZeroMass}{\SI{66}{MeV}}
\providecommand{\VFourEffectiveTrials}{$35.381$}
\providecommand{\VFourSidakValue}{$0.006221$}
\providecommand{\VFourStrongTailMinimum}{$0/300$ at 59, 71, and
72~MeV}
\providecommand{\VFourWeakTailMinimum}{$0/300$ at 21, 65, and
66~MeV}
\providecommand{\VFourTwoSidedTailMinimum}{$2\times0/300$ at
21, 59, 65, 66, 71, and 72~MeV}

\section{Results}
\label{sec:results}

Version 4 reports one internally consistent combined scan: the full 2015 and 2016
samples together with the reviewed 2021 10\% sample. The mass intervals tested for a
signal are unchanged from the recent combined analysis, while every GP is trained on
the wider support declared in Section~\ref{sec:fit-support-policy}. The headline
quantity is the observed 90\% confidence-level upper limit on the shared coupling
$\eps^2$, evaluated with the asymptotic \CLs{} construction. Conditional
background-only pseudoexperiments provide the expected bands and place the observed
limit within the corresponding upper-limit ensemble. Local discovery diagnostics are
reported separately.

\subsection{Declared v4 analysis state}
\label{sec:results-v4-analysis-state}

Table~\ref{tab:results-v4-configuration} records the dataset domains that were fixed
before the observed v4 scan. The search interval controls which resonance hypotheses
enter the result. The fit support controls which spectrum bins may train the GP after
the mass-dependent exclusion is removed. In particular, widening the support does not
add lower- or higher-mass signal hypotheses.

\begin{table}[H]
\centering
\small
{\setlength{\tabcolsep}{5pt}
\begin{tabular}{lccccc}
\toprule
Campaign & Sample & Search [MeV] & GP support [MeV] &
$k_{\min}$ & $k_{\max}$ \\
\midrule
2015 & full & 19--90  & 14--135 & 1.0 & 8 \\
2016 & full & 39--180 & 30--210 & 0.9 & 8 \\
2021 & 10\% & 50--250 & 40--300 & 1.1 & 15 \\
\bottomrule
\end{tabular}
}
\caption{Predeclared v4 data samples, signal-search intervals, and GP fit supports.
The search intervals are identical to those in the preceding combined study. The
wider fit supports restore sideband information on both sides of the endpoint
hypotheses without extending the signal scan. The last two columns give the frozen
lower and upper length-scale factors multiplying the local resolution scale.}
\label{tab:results-v4-configuration}
\end{table}

At a mass covered by more than one campaign, the simultaneous likelihood profiles
one common $\eps^2$. The corresponding signal yields remain dataset dependent through
the luminosity, radiative fraction, mass resolution, efficiency, and branching
inputs described in Section~\ref{sec:event-selection}. The combined curve is therefore
a direct shared-coupling result, not the pointwise envelope of the individual limits.
The minimal visible-dark-photon branching correction above the dimuon threshold is
applied through the same conversion used for the constituent scans.

\subsection{Optimizer-branch review}
\label{sec:results-v4-branch-review}

Every active dataset and mass hypothesis is first evaluated with the frozen v4 card.
A nonfinite result, a visibly discontinuous likelihood branch, or an optimizer state
requiring review is rerun without changing the data, search interval, fit support,
blind geometry, kernel bounds, likelihood, or limit construction. Only the numerical
initialization and optimizer trajectory are allowed to differ. When several repeated
fits are finite, the retained branch is the reviewed solution with the largest
log-marginal likelihood, subject to the same fit-quality checks as the nominal scan.
The repeat seeds, candidate branches, selected branch, and selection reason are kept
in the mass-by-mass optimizer ledger.

No observed limit, expected quantile, or $p$-value is repaired by interpolation. If
unchanged-card repeats do not produce an acceptable branch, the hypothesis remains
explicitly unresolved and is omitted with its status recorded in the ledger. This
replaces the display-level interpolation used for a few v3 review points; none of
those repaired values is carried into the v4 result.

\subsection{Support and kernel-bound audit}
\label{sec:results-v4-support-audit}

The separation between the fit support and the signal-search interval is shown
directly in Figure~\ref{fig:results-v4-support-audit}. At the limiting search
hypotheses, the widened domains leave at least 12 low-side and 138 high-side
training bins for 2015, 20 and 51 for 2016, and 9 and 55 for 2021, after the
mass-dependent exclusion is applied. These are counts of the actual rebinned
training centers rather than continuous-width estimates. This distinction is why
the 2016 lower support is \SI{30}{MeV}: a trial lower edge of \SI{35}{MeV} left no
low-side training center at the \SI{39}{MeV} endpoint with the frozen
\SI{5}{MeV} rebinning.

\begin{figure}[H]
\centering
\graphicorplaceholder{0.92\linewidth}{final_limit_projection_figs/v4_20260803_wide_support/wide_support_search_kernel_audit.pdf}
\caption{Audit of the declared search intervals, wider GP fit supports, and
length-scale-bound occupancy in the reviewed observed scan. The fit support
extends beyond the signal-search interval for all three campaigns. The 2016
length scale reaches its frozen upper bound at all 142 evaluated hypotheses,
whereas the corresponding counts are 0 of 72 for 2015 and 0 of 201 for 2021.
This occupancy is retained as a diagnostic; it is not a reason for a
post-observation retune. Support-matched closure and direct coverage remain
required before treating the conditional ensemble as calibrated.}
\label{fig:results-v4-support-audit}
\end{figure}

All 415 campaign--mass GP states were resolved by unchanged-card fits. Of these,
405 reproduced across all three observed attempts and 10 across two attempts; the
selected ledger contains 135 states from attempt 1, 150 from attempt 2, and 130
from attempt 3. The 2016 boundary occupancy is therefore not an isolated failed
fit or a missing branch. It is a persistent property of the frozen v4 card and is
carried forward visibly rather than removed by changing the kernel bounds after
looking at the observed spectrum.

\subsection{Combined observed limit and conditional expected bands}
\label{sec:results-v4-combined-limit}

Figure~\ref{fig:results-v4-combined-band} is the principal v4 result. The observed
curve uses the asymptotic 90\% \CLs{} upper limit at each reviewed mass hypothesis.
The expected median and central 68\% and 95\% intervals come from 300 background-only
joint pseudoexperiments per mass. A common toy index supplies one pseudo-spectrum for
every campaign active at that mass before the spectra are combined in the
shared-$\eps^2$ likelihood.

\begin{figure}[H]
\centering
\graphicorplaceholder{0.92\linewidth}{final_limit_projection_figs/v4_20260803_wide_support/combined_observed_bands300_minimal_visible.pdf}
\caption{Combined 90\% confidence-level upper limit on $\eps^2$ from the full 2015
and 2016 samples and the 2021 10\% sample. The black curve is the observed asymptotic
\CLs{} result. The expected median and central 68\% and 95\% intervals are obtained
from 300 conditional background-only joint pseudoexperiments at each mass. Search
intervals and fit supports are those in Table~\ref{tab:results-v4-configuration}; no
point is replaced by interpolation. The conversion is the minimal-visible
dark-photon interpretation, including the common dimuon branching correction above
threshold.}
\label{fig:results-v4-combined-band}
\end{figure}

For this ensemble, the reviewed observed-data GP state is held fixed. Each
pseudoexperiment fluctuates the conditional background prediction and event counts,
after which its upper limit is evaluated with the same asymptotic \CLs{} calculation
as the observed spectrum. The GP hyperparameters and sideband fit are not retrained
toy by toy. The bands therefore describe the distribution of limits conditional on
the reviewed GP state. They are not a direct \CLs{} coverage measurement and do not
include the additional variation that would arise from repeating the full training
procedure on every pseudo-spectrum.

With 300 pseudoexperiments, a one-toy change corresponds to an empirical tail-fraction
step of $1/300\simeq0.0033$. The 95\% band is correspondingly less stable than the
central 68\% band because each outer quantile is determined by only the most extreme
few percent of the ensemble. This finite-ensemble precision is shown in the
machine-readable table and retained when the band edges are interpreted.

\subsection{Observed-limit consistency within the pseudoexperiment ensemble}
\label{sec:results-v4-limit-tail}

The strong-, weak-, and two-sided quantities defined in
Section~\ref{sec:methodology} compare the observed upper limit with the same 300
background-only limits used for Figure~\ref{fig:results-v4-combined-band}.
$p_{\mathrm{strong}}$ is the fraction of pseudoexperiment limits at or below the
observed value, while $p_{\mathrm{weak}}$ is the fraction at or above it. The bounded
summary
\[
p_{\mathrm{two}}=
\min\!\left[1,2\min\!\left(p_{\mathrm{strong}},p_{\mathrm{weak}}\right)\right]
\]
flags either an unusually strong or unusually weak observed limit. Figure
\ref{fig:results-v4-limit-tail} presents all three so that the direction of any
departure remains visible.

\begin{figure}[H]
\centering
\graphicorplaceholder{0.92\linewidth}{final_limit_projection_figs/v4_20260803_wide_support/combined_limit_tail_pvalues300.pdf}
\caption{Empirical upper-limit consistency diagnostics from the 300 conditional
background-only joint pseudoexperiments used for the expected bands. The panels show
$p_{\mathrm{strong}}$, $p_{\mathrm{weak}}$, and the bounded
$p_{\mathrm{two}}$. The one-sided values are their inclusive toy counts divided by
300; the two-sided value doubles the smaller one-sided fraction and is capped at
unity. A zero-count tail is reported as 0 of 300, not interpreted as a probability
of zero; any positive floor used to display it on a logarithmic axis is a plotting
convention only. These
quantities locate the observed upper limit within the fixed-GP ensemble and are not
discovery significances or direct coverage tests.}
\label{fig:results-v4-limit-tail}
\end{figure}

In the reviewed table, the smallest strong-tail result is
\VFourStrongTailMinimum, the smallest weak-tail result is
\VFourWeakTailMinimum, and the smallest bounded two-sided result is
\VFourTwoSidedTailMinimum. In each case, the notation means that no
pseudoexperiment entered the indicated empirical tail. It establishes only that the
tail lies below the $1/300$ resolution of this ensemble; it is not a measured
probability of zero. A more granular statement about an extreme tail would require a
larger independent ensemble.

\subsection{Local discovery statistic and analytic trials correction}
\label{sec:results-v4-local-p0}

Figure~\ref{fig:results-v4-local-p0} shows the one-sided fixed-mass asymptotic $p_0$
scan derived from the profiled signal-strength likelihood. The plotted
\v{S}id\'ak curve is the analytic scan-level approximation
\[
p_{\mathrm{Sidak}}(m)
=
1-\left[1-p_0(m)\right]^{N_{\mathrm{eff}}},
\]
using the predeclared effective-trials value for the reviewed mass grid. It is shown
to indicate the size of the look-elsewhere correction under that approximation. It is
not obtained from the 300 expected-limit pseudoexperiments, which contain neither a
full-scan maximum-$q_0$ construction nor a toy-calibrated discovery statistic.

\begin{figure}[H]
\centering
\graphicorplaceholder{0.90\linewidth}{final_limit_projection_figs/v4_20260803_wide_support/combined_local_p0_sidak_reference.pdf}
\caption{Combined fixed-mass asymptotic $p_0$ values and their analytic
\v{S}id\'ak-equivalent scan values. The evaluated hypotheses may be joined to make
the scan legible, but the curve is not smoothed and unresolved hypotheses are not
interpolated. The \v{S}id\'ak conversion is an effective-trials approximation, not a
global $p$-value calibrated with full-scan pseudoexperiments.}
\label{fig:results-v4-local-p0}
\end{figure}

The smallest reviewed local value is \VFourLocalPZeroMinimum{} at
\VFourLocalPZeroMass, corresponding to $Z_{\mathrm{local}}=3.573$. The
resolution-spacing effective-trials value is
\VFourEffectiveTrials, and the corresponding analytic \v{S}id\'ak value is
\VFourSidakValue{} ($Z_{\mathrm{Sidak}}=2.499$). These values come from the final
significance table and remain conceptually separate from the upper-limit ensemble
fractions in Figure~\ref{fig:results-v4-limit-tail}.

\subsection{Matched comparison with the preceding narrow supports}
\label{sec:results-v4-support-comparison}

Figure~\ref{fig:results-v4-wide-narrow-ratio} compares the observed v4 curve with
the preceding narrow-support finalist at matched dataset fractions and mass
hypotheses. The ratio is most visibly displaced in the 2015-only and
2015--2016 overlap regions. Its minimum is 0.2349 at \SI{19}{MeV}, a 76.51\%
tightening relative to the narrow-support result, while its maximum is 2.2567 at
\SI{43}{MeV}, a 125.67\% weakening. Once the scan is supported only by 2021, the
ratio is unity to within approximately $6\times10^{-5}$. The common dimuon
branching correction cancels in this comparison.

\begin{figure}[H]
\centering
\graphicorplaceholder{0.92\linewidth}{final_limit_projection_figs/v4_20260803_wide_support/wide_vs_narrow_observed_limit_ratio.pdf}
\caption{Mass-matched ratio of the v4 observed 90\% \CLs{} limit to the preceding
narrow-support finalist, using the same 2015 full, 2016 full, and 2021 10\% data
fractions. Values below unity are tighter and values above unity are weaker. The
largest changes occur where the widened 2015 or 2016 support contributes; the
2021-only region is unchanged. Because the wider supports were declared before the
observed v4 scan, this is a support-dependence diagnostic rather than an
observed-data optimization or a standalone sensitivity claim.}
\label{fig:results-v4-wide-narrow-ratio}
\end{figure}

The nonmonotonic low-mass response is an important part of the result: extending
the training domain can either strengthen or weaken an observed limit as the fitted
background and its conditional uncertainty change. The comparison therefore does
not justify selecting a support by whichever curve is numerically tighter. Its role
is to document the effect of restoring both sidebands and to motivate the
support-matched closure and coverage checks already identified in
Section~\ref{sec:results-v4-support-audit}.

\subsection{Relation to the v3 review result}
\label{sec:results-v4-v3-context}

Version 3 established the shared-$\eps^2$ likelihood, the 90\% \CLs{} convention,
and the strong-, weak-, and bounded two-sided definitions used here. Its headline
candidate combined the full 2015 sample with only 10\% of 2016 and 1\% of 2021, and
its 2015 and 2016 GP supports ended at the corresponding search boundaries. It also
documented a small number of plot-level interpolations while production reruns were
still outstanding. The v3 tail study used 10,000 pseudoexperiments per mass; v4
retains its definitions but, with the present 300-toy cap, does not inherit that
finer empirical-tail granularity.

Those numerical limits, p-value minima, expected bands, and repaired points are
historical context only. Version 4 refits the full 2016 and 2021 10\% spectra over the
predeclared wider supports, reviews repeated optimizer branches without changing the
card, and recomputes the combined observed scan and its conditional ensemble from
that analysis state. Earlier plots and closure studies remain in
Appendix~\ref{app:prior-validation-results}; they are not mixed with the v4 tables or
used to fill missing v4 hypotheses.
